% =============================================================================
%  全车电气连接原理图 · 图 01 —— 电源与保护（元件级）    v3.0 / 2026-09-20
%  依据：docs/引脚分配表.md v3.0 + 电气连接规范.md §电源
%  决策记录：
%    · 顺序：DC 座 → T3A 慢断保险丝 → P-MOS 反接保护 → 总开关 → 12V_BUS
%      （保险丝在最前，可保护保险丝之后的一切；T3A 是预留值，实车测流后定值）
%    · 反接保护用 AO4407A（SO-8，30V P-MOS）；不用 IRF9540（体积/导通电阻），不优先 AO3401（Vgs 余量）
%    · 3.3V 不由本板产生 —— 核心板自带 RT9193，本板只送 5V（方案 A，2026-09-20 定）
%      ⇒ 本板禁止再把自产 3.3V 灌进核心板 3V3 脚（两颗 LDO 会对打）
%    · 编码器 VCC 实测走 +5V（王政淋 2026-09-20 实测旧板）
%  版面约定（改图请守住）：12V 干线 y=0；5V 母线 y=-3.6；3.3V 母线 y=-7.2；
%    接地干线 y=-10.5 唯一接地符号；12V 支路下拉全落在 x ≤ 18.4，故与下方母线不相交。
%  编译：见同目录 生成接线图.py（在 ASCII 临时目录里跑 xelatex，避开中文路径）
% =============================================================================
\documentclass{article}
\usepackage[paperwidth=55cm,paperheight=25cm,margin=0pt]{geometry}
\usepackage{fontspec}\setmainfont{Latin Modern Roman}
\usepackage{xeCJK}\setCJKmainfont{Microsoft YaHei}
\usepackage{amsmath, amssymb, siunitx}
\usepackage[american]{circuitikz}
\ctikzset{bipoles/length=1.2cm}
\usepackage[active,tightpage]{preview}
\PreviewEnvironment{circuitikz}
\setlength\PreviewBorder{8pt}
\pagestyle{empty}
\begin{document}
\begin{circuitikz}
\small

% ---------------------------------------------------------------- 12V 干线（y = 0）
\draw (0,0) node[ocirc]{} node[above=3pt, align=center, font=\footnotesize]
      {DC005 座\\+ 2P 焊盘\\{\scriptsize 12V\_RAW}}
      to[short] (1.4,0)
      to[fuse, l={$F_1$\ T3A}] (3.4,0)
      to[short] (4.2,0);

% Q1 反接保护（画成模块块，器件已定 AO4407A）
\draw[rounded corners=2pt] (4.2,-0.74) rectangle (6.6,0.74);
\node[align=center, font=\small] at (5.4,0) {$Q_1$\ AO4407A\\P-MOS 反接保护};
\draw (6.6,0) to[switch, l={$SW_1$\ 总开关}] (9.0,0);
\draw (9.0,0) -- (21.4,0);
\fill (9.6,0) circle (1.3pt);                     % 节点：12V_BUS

% Q1 栅极 100k 下拉
\draw (5.4,-0.74) -- (5.4,-1.6) to[R, l=$100\,\mathrm{k}\Omega$] (5.4,-4.8) -- (5.4,-10.5);

% ---- C1 470µF / C2 0.1µF
\draw (10.6,0) -- (10.6,-2.2) to[C] (10.6,-4.2) -- (10.6,-10.5);
\node[anchor=west, font=\footnotesize] at (10.85,-2.75) {$C_1$};
\node[anchor=west, font=\footnotesize] at (10.85,-3.30) {470\,\si{\micro\farad}/25V};
\fill (10.6,0) circle (1.3pt);

\draw (12.6,0) -- (12.6,-2.7) to[C] (12.6,-4.0) -- (12.6,-10.5);
\node[anchor=west, font=\footnotesize] at (12.85,-2.95) {$C_2$};
\node[anchor=west, font=\footnotesize] at (12.85,-3.50) {0.1\,\si{\micro\farad}};
\fill (12.6,0) circle (1.3pt);

% ---- 电池电压采样分压 → PA0
\draw (14.6,0) -- (14.6,-1.6) to[R] (14.6,-3.0);
\node[anchor=east, font=\footnotesize] at (14.35,-1.95) {$R_4$ 100\,k$\Omega$};
\fill (14.6,0) circle (1.3pt);
\fill (14.6,-3.0) circle (1.3pt);
\draw (14.6,-3.0) to[R] (14.6,-5.4) -- (14.6,-10.5);
\node[anchor=east, font=\footnotesize] at (14.35,-3.85) {$R_5$ 10\,k$\Omega$};
\draw (14.6,-3.0) -- (16.4,-3.0) to[C] (16.4,-5.4) -- (16.4,-10.5);
\node[anchor=west, font=\footnotesize] at (16.65,-3.55) {$C_7$ 100\,nF};
\draw (16.4,-3.0) -- (17.6,-3.0) node[ocirc]{};
\node[anchor=south east, font=\footnotesize, align=right] at (17.6,-3.10) {$\to$ PA0\\ADC1\_IN0};

% ---- LED_12V 电源指示灯（不占 IO）
\draw (18.4,0) -- (18.4,-1.6) to[R] (18.4,-2.8) to[leD] (18.4,-4.0) -- (18.4,-10.5);
\node[anchor=west, font=\footnotesize] at (18.65,-1.85) {$R_2$ 1\,k$\Omega$};
\node[anchor=west, font=\footnotesize] at (18.65,-3.05) {$LED_1$ 12V};
\fill (18.4,0) circle (1.3pt);

% ---- VM 输出（大电流，不经过排针）
\draw (19.5,0) -- (19.5,1.6) node[ocirc]{} node[right=2pt, font=\footnotesize]{$\to$ TB6612 \textbf{VM}（大电流，不经过排针）};
\fill (19.5,0) circle (1.3pt);

% ---------------------------------------------------------------- 12V → 5V（MP1584 模块，竖放）
\draw[rounded corners=2pt] (20.4,-0.7) rectangle (22.4,-3.0);
\node[rotate=90, font=\small] at (21.4,-1.85) {MP1584};
\node[anchor=west, font=\footnotesize, align=left] at (22.65,-1.3) {$U_1$ 降压模块\\12V$\to$5V\\{\scriptsize IN+/IN−/OUT+/OUT−}};
\draw (21.4,0) -- (21.4,-0.7);
\draw (21.4,-3.0) -- (21.4,-3.6);

% ---------------------------------------------------------------- 5V 母线（y = -3.6）
\draw (21.4,-3.6) -- (28.0,-3.6);
\fill (21.4,-3.6) circle (1.3pt);

\draw (22.6,-3.6) -- (22.6,-5.0) to[C] (22.6,-6.2) -- (22.6,-10.5);
\node[anchor=west, font=\footnotesize] at (22.85,-5.25) {$C_3$ 10\,\si{\micro\farad}};
\fill (22.6,-3.6) circle (1.3pt);
\draw (24.6,-3.6) -- (24.6,-5.0) to[C] (24.6,-6.2) -- (24.6,-10.5);
\node[anchor=west, font=\footnotesize] at (24.85,-5.25) {$C_4$ 0.1\,\si{\micro\farad}};
\fill (24.6,-3.6) circle (1.3pt);

% ---- 5V 输出（感为手册要求供电稳定，就近再加 100µF + 0.1µF）
\draw (26.4,-3.6) -- (26.4,-2.1) node[ocirc]{} node[right=2pt, font=\footnotesize, align=left]{5V\_BUS：感为灰度 33mA /\\超声波 ×2 / 编码器 ×2};
\fill (26.4,-3.6) circle (1.3pt);

% ---------------------------------------------------------------- 3.3V 由核心板 LDO 供给（本板不产生）
\draw[rounded corners=2pt] (26.8,-4.4) rectangle (29.0,-6.4);
\node[align=center, font=\footnotesize] at (27.9,-5.4) {核心板\\RT9193};
\draw (28.0,-3.6) -- (28.0,-4.4);
\draw (28.0,-6.4) -- (28.0,-7.2);

% ---------------------------------------------------------------- 3.3V 母线（y = -7.2）
\draw (28.0,-7.2) -- (34.6,-7.2);

\draw (29.2,-7.2) -- (29.2,-8.4) to[C] (29.2,-9.2) -- (29.2,-10.5);
\node[anchor=west, font=\footnotesize] at (29.45,-8.55) {$C_5$ 10\,\si{\micro\farad}};
\fill (29.2,-7.2) circle (1.3pt);
\draw (31.2,-7.2) -- (31.2,-8.4) to[C] (31.2,-9.2) -- (31.2,-10.5);
\node[anchor=west, font=\footnotesize] at (31.45,-8.55) {$C_6$ 0.1\,\si{\micro\farad}};
\fill (31.2,-7.2) circle (1.3pt);

\draw (33.0,-7.2) -- (33.0,-5.7) node[ocirc]{} node[right=2pt, font=\footnotesize, align=left]{3V3\_LOGIC：OLED / MPU6050 /\\TB6612 VCC / I2C2 上拉 / 预留（$\leq$300\,mA）};

% ---------------------------------------------------------------- 接地干线（y = -10.5）
\draw (0,-10.5) node[ground]{} -- (31.2,-10.5);
\draw (0,0) -- (0,-10.5);                          % DC 座的地

% ---------------------------------------------------------------- 页脚
\node[anchor=west, font=\footnotesize, align=left] at (0,-11.9)
  {\textbf{硬性要求}：全车只有一块 12V 锂电池（规则硬要求，多电源每轮 $-10\%$）；保险丝 + 反接保护 + 总开关为必需项；\\
   \textbf{核心板 5V 脚与 3V3 脚只能接一个}（本项目接 5V）；本板\textbf{不得}自行产生 3.3V；3V3 负载须核算（RT9193 $\leq$300\,mA）；$F_1$ 为 \textbf{T3A 预留值}，成车实测电流后定值；\\
   12V\_BUS / 5V\_BUS / 3V3\_LOGIC / GND 各留 $\geq$1 个测试点（TP）；12V\_BUS 对地可预留 TVS + 470$\sim$1000\,\si{\micro\farad} 电解 + 100\,nF；\\
   MP1584 首次上电必须\textbf{先脱机}把 OUT 调到 5.0\,V 再接后级；开关节点/电感附近不走 I2C、编码器、蓝牙天线等敏感信号；\\
   每个模块接口的 5V 支路串 0$\Omega$ 跳线或自恢复保险丝，挑故障时可单独断开。};

\end{circuitikz}
\end{document}
